%% kapa: exam=1 round=37 year=2026 period=1 form=A subject=economics subject_ko=경제학원론 seq=2 number=42
%% source: Q-Net 감정평가사 기출문제 (cst00309) artlSeq=5258924 file=raw/1008/5258924/2026년 제37회 감정평가사 1차 1교시.pdf
%% conversion: kordoc 4.24.51 -> tools/gen1.py (auto); review=proofread-sonnet; flags=none
\begin{question}{42}
A국의 휘발유 시장은 완전경쟁이며, 수요함수는 $Q_{D}=1,000-\frac{1}{2}P$이고, 공급함수는 $Q_{S}=\frac{1}{3}P-\frac{500}{3}$이다. 정부가 휘발유 1리터당 500원의 종량세를 공급자에게 부과할 때, 다음 설명 중 옳은 것을 모두 고른 것은? (단, $Q$는 수량(리터), $P$는 가격)
\begin{framedtext}
ㄱ. 세금은 모두 공급자가 부담한다. \par
ㄴ. 균형량은 종량세 부과 전에 비해 감소한다. \par
ㄷ. 소비자 잉여는 종량세 부과 전에 비해 증가한다. \par
ㄹ. 동일한 종량세를 수요자에게 부과해도 세후 균형량은 동일하다. \par
\end{framedtext}
\begin{choices}
  \item ㄱ, ㄴ
  \item ㄱ, ㄷ
  \item ㄴ, ㄷ
  \item ㄴ, ㄹ
  \item ㄷ, ㄹ
\end{choices}
\end{question}
\answer{4}
